\begin{frame}{Un tema recurrente en matemáticas}
    \vfill
       \center{\LARGE{\alert{Local/Finito vs Global/Infinito}}}
    \vfill
\end{frame}
\begin{frame}{Compacidad}
   En topología:
   \begin{Def}
      Un espacio topológico $X$ se dice compacto si todo cubrimiento por abiertos de $X$ tiene un subcubrimiento finito.\vspace{0.3cm}
   \end{Def}
   \vspace{0.3cm}
   \onslide<2->
   En lógica:
   \begin{Tma}[De compacidad]
      Sean $L$ un lenguaje de primer orden y $\mathcal{E}$ un conjunto de $L$-sentencias.
      \center{Si $\mathcal{E}$ es finitamente satisfacible entonces $\mathcal{E}$ es satisfacible.\vspace{0.3cm}}
   \end{Tma}
\end{frame}
\begin{frame}{Compacidad}
    \begin{Tma}[De compacidad]
      Sean $L$ un lenguaje de primer orden y $\mathcal{E}$ un conjunto de $L$-sentencias.
      \center{Si $\mathcal{E}$ es finitamente satisfacible entonces $\mathcal{E}$ es satisfacible.\vspace{0.3cm}}
   \end{Tma}
   \vspace{0.3cm}
   \begin{itemize}
      \item $\mathcal{E}$ se dice finitamente satisfacible si todo subconjunto finito de $\mathcal{E}$ tiene un modelo.
      \item<2-> $\mathcal{E}$ se dice satisfacible si tiene un modelo.
   \end{itemize}
\end{frame}


